\documentclass[11pt,a4paper]{article}
\usepackage[utf8]{inputenc}
\usepackage[vietnamese]{babel}
\usepackage{amsmath,amssymb}
\usepackage{geometry}
\geometry{a4paper, total={180mm,260mm}, left=15mm, right=15mm, top=18mm, bottom=18mm}
\usepackage{tikz}
\usetikzlibrary{shapes.geometric, arrows.meta, positioning, calc}
\usepackage{tcolorbox}
\tcbuselibrary{skins,breakable}
\usepackage{xcolor}
\usepackage{multicol}
\usepackage{array}
\usepackage{fontawesome}
\usepackage{fancyhdr}

\pagestyle{fancy}
\fancyhf{}
\renewcommand{\headrulewidth}{0.5pt}
\renewcommand{\footrulewidth}{0.5pt}
\fancyhead[L]{\small\bfseries\color{teal}\faBook\ SỔ TAY KIẾN THỨC TOÁN 3}
\fancyhead[R]{\small\bfseries\color{teal}Lớp Toán Thầy Lĩnh -- loptoanthaylinh.vn}
\fancyfoot[C]{\small\color{gray}-- Sổ tay kiến thức: Bài 76. Ôn tập các số trong phạm vi 10 000, 100 000 --}

\definecolor{primary}{HTML}{005F73}
\definecolor{secondary}{HTML}{0A9396}
\definecolor{accent}{HTML}{EE9B00}
\definecolor{danger}{HTML}{AE2012}
\definecolor{boxbg}{HTML}{F4F6F8}

\begin{document}

% --- BANNER TIÊU ĐỀ ---
\begin{tcolorbox}[
  enhanced,
  colback=white,
  colframe=primary,
  arc=3mm,
  boxrule=1.2pt,
  left=4mm, right=4mm, top=2.5mm, bottom=2.5mm,
  drop shadow=black!10!white
]
  \centering
  {\Large\bfseries\color{primary}\faBookmark\ BÀI 76: ÔN TẬP CÁC SỐ PHẠM VI 10 000, 100 000}\\[1.5mm]
  {\small\itshape\color{black!75}Ôn tập cuối năm: Đọc viết cấu tạo số, so sánh dãy số, số La Mã và làm tròn số 4 cấp độ trên 1 trang A4}
\end{tcolorbox}

\vspace{-2mm}


\noindent
\begin{minipage}[t]{0.49\textwidth}\vspace{0pt}
\begin{tcolorbox}[
  enhanced,
  colback=blue!4!white,
  colframe=blue!75!black,
  arc=2mm, boxrule=0.8pt,
  left=2mm, right=2mm, top=1.8mm, bottom=1.8mm,
  title={\footnotesize\bfseries\color{white}\faKey\ 1. Cấu tạo số và số La Mã}
]
{\footnotesize
\begin{itemize}\setlength{\itemsep}{1.5pt}\setlength{\leftskip}{-2.5mm}
  \item Số có 5 chữ số gồm 5 hàng: Chục nghìn, Nghìn, Trăm, Chục, Đơn vị.
  \item Số $100\,000$ (Một trăm nghìn): Số nhỏ nhất có 6 chữ số.
  \item \textbf{Chữ số La Mã cốt lõi:}\\
  $\text{I}=1, \text{V}=5, \text{X}=10$; $\text{IV}=4, \text{VI}=6, \text{IX}=9, \text{XI}=11, \text{XIX}=19, \text{XX}=20$.
\end{itemize}
}
\end{tcolorbox}
\end{minipage}
\hfill
\begin{minipage}[t]{0.49\textwidth}\vspace{0pt}
\begin{tcolorbox}[
  enhanced,
  colback=teal!4!white,
  colframe=teal!75!black,
  arc=2mm, boxrule=0.8pt,
  left=2mm, right=2mm, top=1.8mm, bottom=1.8mm,
  title={\footnotesize\bfseries\color{white}\faSliders\ 2. So sánh và Làm tròn số}
]
{\footnotesize
\begin{itemize}\setlength{\itemsep}{1.5pt}\setlength{\leftskip}{-2.5mm}
  \item \textbf{So sánh:} Số nhiều chữ số hơn thì lớn hơn; cùng số chữ số thì so sánh từ trái sang phải.
  \item \textbf{Quy tắc mốc 5 làm tròn:}\\
  -- Tròn chục $\to$ xét hàng đơn vị.\\
  -- Tròn trăm $\to$ xét hàng chục.\\
  -- Tròn nghìn $\to$ xét hàng trăm.\\
  -- Tròn chục nghìn $\to$ xét hàng nghìn.
\end{itemize}
}
\end{tcolorbox}
\end{minipage}

\vspace{-2.5mm}
% --- KHỐI 3: SƠ ĐỒ TƯ DUY TIKZ ---
\noindent
\begin{tcolorbox}[
  enhanced,
  colback=orange!4!white,
  colframe=orange!80!black,
  arc=2mm, boxrule=0.8pt,
  left=2mm, right=2mm, top=1.5mm, bottom=1.5mm,
  title={\footnotesize\bfseries\color{white}\faSitemap\ SƠ ĐỒ TƯ DUY: HỆ THỐNG KIẾN THỨC SỐ HỌC TOÁN 3}
]
\centering
\begin{tikzpicture}[
  box/.style={draw, rounded corners=2pt, align=center, font=\scriptsize\bfseries, inner sep=2.5pt, minimum height=0.75cm},
  arr/.style={-{Stealth[length=2mm]}, thick}
]
  \node[box, fill=purple!15, draw=purple!80!black, text width=2.6cm] (h1) at (0, 0) {CẤU TẠO SỐ\\Nghìn, Chục nghìn\\$100\,000$ (6 chữ số)};
  \node[box, fill=teal!15, draw=teal!80!black, text width=2.6cm] (h2) at (3.5, 0) {SO SÁNH THỨ TỰ\\Xét từ trái sang phải\\Tăng dần / Giảm dần};
  \node[box, fill=blue!12, draw=blue!80!black, text width=2.6cm] (h3) at (7.0, 0) {CHỮ SỐ LA MÃ\\$\text{I}, \text{V}, \text{X}$\\Đọc đồng hồ La Mã};
  \node[box, fill=orange!15, draw=orange!80!black, text width=2.6cm] (h4) at (10.5, 0) {LÀM TRÒN SỐ\\Mốc 5: Xuống / Lên\\4 cấp độ làm tròn};

  \draw[arr, purple!80!black] (h1) -- (h2);
  \draw[arr, teal!80!black] (h2) -- (h3);
  \draw[arr, blue!80!black] (h3) -- (h4);
\end{tikzpicture}
\end{tcolorbox}

\vspace{-2.5mm}
% --- KHỐI 4: VÍ DỤ & MẸO TRÁNH BẪY ---
\noindent
\begin{minipage}[t]{0.49\textwidth}\vspace{0pt}
\begin{tcolorbox}[
  enhanced,
  colback=green!4!white,
  colframe=green!60!black,
  arc=2mm, boxrule=0.8pt,
  left=2mm, right=2mm, top=1.8mm, bottom=1.8mm,
  title={\footnotesize\bfseries\color{white}\faCheckCircle\ Ví dụ thực hành chuẩn mực}
]
{\footnotesize
\textbf{1. Viết số La Mã:}\\
Số $14$ viết là: $\mathbf{\text{XIV}}$.\\
Số $19$ viết là: $\mathbf{\text{XIX}}$.\\
\textbf{2. Làm tròn phối hợp:}\\
Cho số $47\,528$:\\
-- Làm tròn đến hàng trăm: $\mathbf{47\,500}$ (chục là $2 < 5$).\\
-- Làm tròn đến hàng nghìn: $\mathbf{48\,000}$ (trăm là $5 \ge 5$).\\
-- Làm tròn đến hàng chục nghìn: $\mathbf{50\,000}$ (nghìn là $7 \ge 5$).
}
\end{tcolorbox}
\end{minipage}
\hfill
\begin{minipage}[t]{0.49\textwidth}\vspace{0pt}
\begin{tcolorbox}[
  enhanced,
  colback=red!4!white,
  colframe=red!70!black,
  arc=2mm, boxrule=0.8pt,
  left=2mm, right=2mm, top=1.8mm, bottom=1.8mm,
  title={\footnotesize\bfseries\color{white}\faExclamationTriangle\ Mẹo tránh bẫy trắc nghiệm}
]
{\footnotesize
\begin{itemize}\setlength{\itemsep}{1.5pt}\setlength{\leftskip}{-2.5mm}
  \item \textbf{Bẫy làm tròn dây chuyền:}\\
  Làm tròn đến hàng chục nghìn của số $47\,528$ thì chỉ nhìn vào chữ số hàng nghìn ($7 \ge 5$), không được làm tròn từng bước từ đơn vị lên!
  \item Không nhầm $\text{IV}$ ($4$) với $\text{VI}$ ($6$).
\end{itemize}
}
\end{tcolorbox}
\end{minipage}


\end{document}
